\bookchapter{Глава 7: Полный разбор передовых функций}

\chapterintro{Технологии не стоят на месте, и функции телефонов постоянно развиваются. Спутниковая связь, связь без сети, обнаружение аварий, технология SparkLink… Эти функции, звучащие как научная фантастика, уже вошли в нашу жизнь. Давайте разберёмся в этих передовых технологиях.}

\booksection{Спутниковая связь: помощь даже без сигнала}

Спутниковая связь — один из крупнейших прорывов в мобильной сфере последних лет. Когда вы в глуши, в горах, и телефон полностью лишён сигнала, спутниковая связь может стать спасительной соломинкой.

\booksubsection{Принцип работы:}

Телефон напрямую связывается с коммуникационным спутником на орбите и отправляет сигнал бедствия. В настоящее время поддерживается в основном отправка SMS и координат местоположения, звонки недоступны.

\booksubsection{Поддерживаемые модели:}

\begin{itemize}
  \item \textbf{Apple}: серия iPhone 14 и новее, поддерживает экстренный SOS через спутник. Серия iPhone 17 дополнительно поддерживает отправку iMessage и SMS через спутник, а также расширяет покрытие на большее количество регионов.
  \item \textbf{Huawei}: начиная с серии Mate 50 поддерживает спутниковые сообщения Beidou (только текст), серия Mate 60 Pro дополнительно поддерживает спутниковые звонки Tiantong (можно звонить), возможности спутниковой связи в серии Mate 70 дополнительно усилены.
  \item \textbf{Xiaomi/OPPO/vivo}: с 2024 года флагманские модели последовательно поддерживают функцию спутниковых сообщений Beidou.
  \item \textbf{Тенденции развития}: спутниковая связь распространяется от эксклюзивной функции премиум-флагманов вниз, в будущем её могут получить и модели среднего класса. Китайская спутниковая система Beidou имеет наиболее широкую поддержку среди отечественных телефонов.
\end{itemize}

\booksubsection{Сценарии использования:}

\begin{itemize}
  \item Походы, альпинизм, чрезвычайные ситуации на природе.
  \item Чрезвычайные ситуации в море.
  \item Нарушение наземной связи из-за стихийных бедствий.
\end{itemize}

\booksubsection{Способ использования:}

\begin{enumerate}
  \item Убедитесь, что вы на открытом пространстве с видом на небо.
  \item Откройте функцию спутниковой связи телефона (обычно в интерфейсе экстренного вызова).
  \item Следуйте инструкциям на экране и направьте телефон в сторону спутника.
  \item Отправьте предустановленное сообщение о помощи и координаты.
\end{enumerate}

\begin{tcbinfo}{Маленькая заметка}
Для спутниковой связи требуется специальная аппаратная антенна в телефоне, не все телефоны её поддерживают. Кроме того, спутниковая связь обычно требует платной подписки, но функция экстренного вызова, как правило, бесплатна.

\end{tcbinfo}

\booksection{Связь без сети: новый способ ближней коммуникации}

«Связь без сети» — это технология прямой связи между двумя телефонами при отсутствии мобильной сети и Wi-Fi.

\booksubsection{Технический принцип:}

Использование Bluetooth или выделенных радиочастот для установления прямого соединения между устройствами на коротком расстоянии (обычно от нескольких сотен метров до нескольких километров).

\booksubsection{Сценарии применения:}

\begin{itemize}
  \item \textbf{Активный отдых}: связь между участниками похода, кемпинга.
  \item \textbf{Спасательные операции}: при нарушении связи после землетрясений и наводнений.
  \item \textbf{Массовые мероприятия}: концерты, спортивные события, где много людей и плохой сигнал.
\end{itemize}

\booksubsection{Поддерживаемые платформы:}

\begin{itemize}
  \item \textbf{Apple}: функции NameDrop и «Поиск друзей офлайн» в iOS 17.
  \item \textbf{Huawei}: приложение ChangLian поддерживает звонки без сети.
  \item \textbf{Xiaomi}: некоторые функции центра Xiaomi Miaoxiang.
\end{itemize}

\booksection{Обнаружение аварий: спасение жизни в критический момент}

Обнаружение аварий — это функция, использующая датчики телефона для выявления серьёзных ДТП и автоматического вызова экстренных служб.

\booksubsection{Принцип работы:}

Телефон с помощью акселерометра, гироскопа, барометра, микрофона и других датчиков комплексно определяет, произошло ли серьёзное столкновение. При обнаружении аварии телефон:

\begin{enumerate}
  \item Показывает экран обратного отсчёта с вопросом, всё ли в порядке.
  \item Если вы не ответили до окончания отсчёта, автоматически набирает экстренный номер.
  \item Отправляет ваше местоположение экстренным контактам.
\end{enumerate}

\booksubsection{Поддерживаемые модели:}

\begin{itemize}
  \item \textbf{Apple}: серия iPhone 14 и новее.
  \item \textbf{Android}: некоторые премиум-модели постепенно получают поддержку.
\end{itemize}

\booksubsection{На заметку:}

\begin{itemize}
  \item Функция должна быть включена для использования.
  \item Возможны ложные срабатывания (например, на американских горках или при интенсивных упражнениях) — просто отмените вовремя.
  \item Зависит от датчиков и алгоритмов, не на 100\% надёжна, но это дополнительный уровень защиты.
\end{itemize}

\booksection{Технология SparkLink: беспроводная связь следующего поколения}

SparkLink — это технология беспроводной связи нового поколения, разработанная под руководством Huawei, цель которой — заменить Bluetooth и Wi-Fi в определённых сценариях. По ключевым показателям она сопоставима с существующими решениями или превосходит их: скорость передачи по радиоканалу до 900 Мбит/с (значительно выше 50 Мбит/с у Bluetooth), задержка всего 20 мкс (у Bluetooth около 15 мс), одна сеть поддерживает до 4096 одновременно подключённых устройств (у Bluetooth обычно до 8).

\booksubsection{Технические особенности:}

\begin{itemize}
  \item \textbf{Сверхнизкая задержка}: сквозная задержка всего 20 микросекунд (0,02 мс), на два порядка ниже, чем у Bluetooth, практически незаметна в играх и синхронизации музыки.
  \item \textbf{Высокая скорость}: пиковая скорость одного канала 900 Мбит/с, в десятки раз выше, чем у Bluetooth Audio, позволяет передавать lossless-аудио и даже HD-видео.
  \item \textbf{Низкое энергопотребление}: энергопотребление в режиме ожидания сопоставимо с Bluetooth, но эффективность передачи выше, расход энергии на единицу данных ниже.
  \item \textbf{Широкое покрытие + высокая плотность}: дальность соединения более чем на 50\% больше, чем у Bluetooth, одна сеть поддерживает тысячи устройств одновременно.
\end{itemize}

\booksubsection{Сценарии применения:}

\begin{itemize}
  \item \textbf{Беспроводные наушники}: меньшая задержка, синхронизация звука и видео в играх.
  \item \textbf{Умные часы}: более стабильное соединение.
  \item \textbf{Умный дом}: больше устройств одновременно на связи.
  \item \textbf{Автомобильная интеграция}: более плавное соединение телефона с автомобилем.
\end{itemize}

\booksubsection{Поддерживаемые устройства:}

В настоящее время в основном устройства экосистемы Huawei, такие как серия Huawei Mate 60, некоторые наушники и часы Huawei. В будущем ожидается поддержка от большего числа брендов.

\textbf{Перспективы технологии}\par

Технология SparkLink представляет собой будущее направление беспроводной связи. С присоединением новых производителей мы увидим больше устройств с поддержкой SparkLink, обеспечивающих превосходный беспроводной опыт.

\begin{tcbpractice}{Проверьте: поддерживает ли ваш телефон эти функции}
\begin{enumerate}
  \item Поищите «Спутник» в настройках — посмотрите, есть ли опции спутниковой связи.
  \item Поищите «Обнаружение аварий» или «Экстренный SOS» — проверьте, включено ли.
  \item Если ваш телефон Huawei, посмотрите, поддерживает ли он подключение устройств SparkLink.
  \item Узнайте, как пользоваться этими функциями — в критический момент они могут пригодиться.
\end{enumerate}

\end{tcbpractice}
